\subsection{Reduced Density Matrix}\label{sec:density-matrix-reduced}

\valert{
Information about the many-body entanglement of a state is encoded in the reduced density matrix $\rho_A$, knowing which is the first step in investigating the entanglement structure. Here, we introduce an intuitive approach to obtain the reduced density matrix, where the identity between the spin representation and the Majorana-fermion representation is demonstrated again, similar to the previous section. Nevertheless, we also provide a rigorous derivation of how to carry out the trace directly in Appendix~\ref{appendix:reduced-density0matrix}. 
}

The property of a reduced density matrix is pretty much the same as the density matrix of the entire system, since both are designed to give the expectation values of observables in the way that $\anglelr{\mathcal{O}_A}=\Tr\lr{\rho\mathcal{O}_A}=\Tr_A\lr{\rho_A\mathcal{O}_A}$, for any observable $\mathcal{O}_A$ with support $A$. In the Kitaev model, observables possessing finite expectation values are string operators. Therefore, the reduced density matrix can be expressed in terms of string operators $\Sigma_\mathfrak{D}$, with $\mathfrak{D}$ entirely inside $A$, accompanied by their expectation values $\mathfrak{C}_\mathfrak{D}$,
\begin{equation}\label{equ.rho-A-raw-guess}
    \rho_A \propto\sum_{\mathfrak{D}\subset A}\mathfrak{C}_\mathfrak{D} \Sigma_\mathfrak{D}
\end{equation}
The equivalent configurations can be sorted in the same way as in the previous section, then the summation can be rewritten as a summing of equivalent classes with the projector $\mathcal{P}^{(p)}_A$ of dimension $2^{N^A-N^A_p}$, which only includes plaquettes inside $A$,
\begin{equation}\label{equ.p-proj-A}
    \mathcal{P}^{(p)}_A = \prod_{p\in A}\lr{\frac{\mathds{1}+W_p}{2}}
\end{equation}
The final form of the reduced density matrix is then
\[
\rho_A = C\mathcal{P}^{(p)}_A\sum_{[\mathfrak{D}]\subset A}\mathfrak{C}_{\mathfrak{D}}\Sigma_{\mathfrak{D}}
\]
The condition $[\mathfrak{D}]\subset A$ below the summation means choosing equivalence classes such that each class contains at least one configuration that is entirely contained within $A$. Recall the correspondence between endpoint sets and equivalent classes, the condition is also equivalent to choosing classes $[\mathfrak{D}]$ such that $\mrm{EP}(\mathfrak{D})\subset A$. The constant $C$ is a normalization coefficient to be determined by imposing the normalization condition $\Tr_A\rho_A=1$,
\begin{align*}
    \Tr_A\rho_A &= C\sum_{[\mathfrak{D}]\subset A} \mathfrak{C}_\mathfrak{D} \Tr_A\lrm{\mathcal{P}^{(p)}_A \Sigma_\mathfrak{D}} \\
    &= \frac{C}{2^{N^A_p}}\sum_{\mathcal{A}_p}\sum_{[\mathfrak{D}]\subset A} \mathfrak{C}_\mathfrak{D} \Tr_A\lrm{ \lr{\prod_{p\in\mathcal{A}_p}W_p} \Sigma_\mathfrak{D}} \\
    &= \frac{C}{2^{N^A_p}}\Tr_A\mathds{1}_A = C2^{N^A - N^A_p} = 1
\end{align*}
Therefore, the final form of the reduced density matrix is as follows.
\begin{equation}\label{equ.red-den-mat-sim-con}
    \rho_A = \frac{1}{2^{N^A-N^A_p}} \mathcal{P}^{(p)}_A\sum_{[\mathfrak{D}]\subset A}\mathfrak{C}_{\mathfrak{D}}\Sigma_{\mathfrak{D}}
\end{equation}
The form \eqref{equ.red-den-mat-sim-con} is the same as the density matrix \eqref{equ.density-matrix} as we anticipated. 

